\documentclass[main.tex]{subfiles}
\pagenumbering{arabic}
\newcommand\chapterlabel{ChMIX-biomoleculesinsolution}\setcounter{figurenewcounter}{0}\setcounter{tablenewcounter}{0}\setcounter{formulanewcounter}{0}
\setcounter{chapter}{5}

\begin{document}



\chapter[atoms  ]{The atom }
\renewcommand\chapterlabel{Ch-Table}
 \begin{marginfigure}\includegraphics{../\chapterlabel/figure1} \end{marginfigure}\import{../\chapterlabel/files/}{ChapterIntro}

\begin{marginfigure}%LEARNING GOALS BOX
\begin{mytcbox}{GOALS}
\renewcommand\chapterlabel{Ch-Table}
\begin{center}\textcolor{blue}{\faEnvira Chemistry and molecules }\end{center}
\begin{enumerate}[label=\protect\circled{\color{white}\arabic*}]
  \import{../\chapterlabel/files/}{SectionGoal-The-periodic-table}
  \import{../\chapterlabel/files/}{SectionGoal-The-atom}
\end{enumerate}

\renewcommand\chapterlabel{Ch-acidbase}
\begin{center}\textcolor{blue}{\faFlask Acids and bases }\end{center}
\begin{enumerate}[label=\protect\circled{\color{white}\arabic*}]
\import{../\chapterlabel/files/}{SectionGoal-Titrations}
\end{enumerate}

\renewcommand\chapterlabel{Ch-radiation}
\begin{center}\textcolor{blue}{\faTint  Electrolytes and solutions }\end{center}
\begin{enumerate}[label=\protect\circled{\color{white}\arabic*}]
\import{../\chapterlabel/files/}{SectionGoal-Electronic-configuration-Full-electron-Configuration}
\import{../\chapterlabel/files/}{SectionGoal-Periodic-Trends}
\end{enumerate}

\renewcommand\chapterlabel{Ch-mole}
\begin{center}\textcolor{blue}{\faFire Reactions and energy} \end{center}
\begin{enumerate}[label=\protect\circled{\color{white}\arabic*}]
 \import{../\chapterlabel/files/}{SectionGoal-Stoichiometry-and-mass-calculations-Mole-Mole-ratio}
\end{enumerate}

\renewcommand\chapterlabel{Ch-nuclear}
\begin{center}\textcolor{blue}{\faFlash The chemistry of cancer} \end{center}
\begin{enumerate}[label=\protect\circled{\color{white}\arabic*}]
\import{../Ch-nuclear/../\chapterlabel/files/}{SectionGoal-Unknown-isotopes-in-nuclear-reactions}
\end{enumerate}
	
%\begin{center}\textcolor{blue}{\faFlash Let's calculate} \end{center}
%\begin{enumerate}[label=\protect\circled{\color{white}\arabic*}]
%\import{../Ch-measurements/../\chapterlabel/files/}{SectionGoal-Using-Conversion-Factors-Square-or-cubic-units}
%\import{../Ch-measurements/../\chapterlabel/files/}{SectionGoal-Using-Conversion-Factors-Units-of-volume}
%\end{enumerate}
%			 				 




\end{mytcbox}
\vspace{1cm}
%\begin{tcolorbox}[enhanced,colback=red!5!white,colframe=black!50!red,boxrule=1pt,
%  arc=0pt,outer arc=0pt,drop heavy lifted shadow]
%\faGears\ 
% \import{../Ch-electrolytes/../\chapterlabel/files/}{ChapterDiscussion}
% \end{tcolorbox}
\end{marginfigure}%LEARNING GOALS BOX






\renewcommand\chapterlabel{Ch-Table}
\section{\faEnvira The periodic table}\import{../\chapterlabel/files/}{SectionIntro-The-periodic-table}
\sloppy\begin{description}
\import{../\chapterlabel/files/}{SubSection-The-periodic-table-Elements-and-Symbols}
\import{../\chapterlabel/problems/}{SampleProblem9}
\import{../\chapterlabel/files/}{SubSection-The-periodic-table-Periods-and-groups}
\import{../\chapterlabel/files/}{SubSection-The-periodic-table-Properties-in-the-periodic-table}
\import{../\chapterlabel/files/}{SubSection-The-periodic-table-Metals-Nonmetals-and-Metalloids}
\import{../\chapterlabel/problems/}{SampleProblem1}
\import{../\chapterlabel/files/}{SubSection-The-periodic-table-Classification-of-elements-in-terms-of-groups}
\import{../\chapterlabel/files/}{SubSection-The-periodic-table-How-to-classify-Hydrogen}
 \hspace{0cm}\import{../\chapterlabel/files/}{Figure-The-periodic-table}
 \end{description}
 
\section{\faEnvira The atom}
\import{../\chapterlabel/files/}{SectionIntro-The-atom}
\sloppy\begin{description}
\import{../\chapterlabel/files/}{SubSection-The-atom-Atomic-Structure}
\import{../\chapterlabel/files/}{SubSection-The-atom-Atomic-and-mass-number}
\import{../\chapterlabel/problems/}{SampleProblem2}
\import{../\chapterlabel/files/}{SubSection-The-atom-Isotopes}
\import{../\chapterlabel/files/}{SubSection-The-atom-Average-atomic-mass}
\import{../\chapterlabel/problems/}{SampleProblem3}
 \import{../\chapterlabel/files/}{Table-Isotopes}
\end{description}






\renewcommand\chapterlabel{Ch-radiation}
\section{\faTint Electronic configuration of an atom}\import{../\chapterlabel/files/}{SectionIntro-Electronic-configuration-of-an-atom}
\sloppy\begin{description}
\item[\docfilehook{Electron energy levels and sublevels}{}] \import{../\chapterlabel/files/}{SubSection-Electronic-configuration-Electron-energy-levels-and-sublevels}
  \import{../\chapterlabel/problems/}{SampleProblem8}
\item[\docfilehook{The Pauli exclusion principle}{}] \import{../\chapterlabel/files/}{SubSection-Electronic-configuration-The-Pauli-exclusion-principle}
\stepcounter{figurenewcounter}   \refstepcounter{figure}  \label{Fig:{\chapterlabel}\thefigurenewcounter} 
\item[\docfilehook{Orbital Filling: the aufbau principle}{}]\import{../\chapterlabel/files/}{SubSection-Electronic-configuration-Orbital-Filling-the-aufbau-principle}
\import{../\chapterlabel/files/}{SideFigure-Orbital-filling-diagram}
 \item[\docfilehook{Full electron Configuration}{}] \import{../\chapterlabel/files/}{SubSection-Electronic-configuration-Full-electron-Configuration}
  \import{../\chapterlabel/problems/}{SampleProblem9}
\item[\docfilehook{Abbreviated Electron Configuration and orbital diagrams}{}] 
\import{../\chapterlabel/files/}{SubSection-Electronic-configuration-Abbreviated-Electron-Configuration-and-orbital-diagrams}
  \import{../\chapterlabel/problems/}{SampleProblem10}
  \item[\docfilehook{Hund's rule}{}] \import{../\chapterlabel/files/}{SubSection-Electronic-configuration-Hund's-rule}
\import{../\chapterlabel/problems/}{SampleProblem12}
\end{description}

\newpage
\section{\faTint  Periodic Trends}\import{../\chapterlabel/files/}{SectionIntro-Periodic-Trends}
\sloppy\begin{description}
\item[\docfilehook{Effective charge}{}] \import{../\chapterlabel/files/}{SubSection-Periodic-Trends-Effective-charge}
   \import{../\chapterlabel/files/}{Figure-Atomic-radius}
\import{../\chapterlabel/files/}{Figure-Atomic-radius-plot}
\item[\docfilehook{Atomic radii}{}] \import{../\chapterlabel/files/}{SubSection-Periodic-Trends-Atomic-radius}
\begin{fullwidth}\import{../\chapterlabel/files/}{Table-Ionization-energies}\end{fullwidth}\newpage
%\import{../\chapterlabel/files/}{Figure-Ionization-energy-plot}
\item[\docfilehook{Ionization Energy}{}] \import{../\chapterlabel/files/}{SubSection-Periodic-Trends-Ionization-Energy}
\item[\docfilehook{Electron affinity}{}] \import{../\chapterlabel/files/}{SubSection-Periodic-Trends-Electron-affinity}
\hspace{-6cm}  \import{../\chapterlabel/files/}{Figure-Periodic-propeties}
\item[\docfilehook{Electronegativity, EN}{}] \import{../\chapterlabel/files/}{SubSection-Periodic-Trends-Electronegativity,-EN}
\item[\docfilehook{Metallic character }{Metallic character }]\import{../\chapterlabel/files/}{SubSection-Periodic-Trends-Metallic-character-}
  \import{../\chapterlabel/problems/}{SampleProblem11}
\item[\docfilehook{Ionic radius}{}]  \import{../\chapterlabel/files/}{SubSection-Periodic-Trends-Ionic-radius}
\end{description}

\renewcommand\chapterlabel{Ch-acidbase}
\section{\faFlask Titrations}\import{../\chapterlabel/files/}{SectionIntro-Titrations}
\sloppy \begin{description}
\item[\docfilehook{ Neutralization Reactions}{}] \import{../\chapterlabel/files/}{SubSection-Titrations-Neutralization-Reactions}
\item[\docfilehook{ Endpoint formula}{}] \import{../\chapterlabel/files/}{SubSection-Titrations-Endpoint-formula}
  \import{../\chapterlabel/problems/}{SampleProblem17}
\end{description}




\renewcommand\chapterlabel{Ch-nuclear}
\section{\faFlash Unknown isotopes in nuclear reactions}
\import{../\chapterlabel/files/}{SectionIntro-Unknown-isotopes-in-nuclear-reactions}
  \import{../\chapterlabel/problems/}{SampleProblem3}




\renewcommand\chapterlabel{Ch-mole}
\section{\faFire Stoichiometry }
\import{../\chapterlabel/files/}{SectionIntro-Stoichiometry-and-mass-calculations}
\sloppy \begin{description}
\import{../\chapterlabel/files/}{SubSection-Stoichiometry-and-mass-calculations-Mole-Mole-ratio}
\import{../\chapterlabel/files/}{SubSection-Stoichiometry-and-mass-calculations-Reactants-to-products}
\import{../\chapterlabel/files/}{SubSection-Stoichiometry-and-mass-calculations-Reactant-to-a-different-reactant}
\import{../\chapterlabel/problems/}{SampleProblem6}
\end{description}















 






\end{document}
 
